\documentclass{beamer}

\usepackage[version=4]{mhchem}
\usepackage{chemfig}

\usepackage{hyperref}
\usepackage{graphicx}
\usepackage[T1]{fontenc}
\usepackage[dvipsnames]{xcolor}

\usetheme{moloch}
\setbeamertemplate{page number in head/foot}[totalframenumber]
\setbeamertemplate{section in toc}[sections numbered]
\setbeamertemplate{section page}[progressbar]
\setbeamertemplate{title}[plain]

\setbeamercolor{palette primary}{bg=
Maroon, fg=Cerulean!15!white}
\setbeamercolor{palette secondary}{fg=blue,bg=red}
\setbeamercolor{background canvas}{bg=GreenYellow!12!white}
\setbeamercolor{normal text}{fg=RedViolet!10!black}
\setbeamercolor{frametitle}{
fg=Yellow!15!white}
\setbeamercolor{title}{fg=MidnightBlue}
\usebeamercolor[fg]{normal text}

%wg dokumentacji to należy dodać expl3, amsmath and calc
%\usepackage{expl3}
\usepackage{amsmath}
%\usepackage{calc}

\title{\centerline{\color{Maroon}tryb chemiczny -- mchem i chemfig}}
\subtitle{\centerline{warsztaty \LaTeX \ 2025}}
\author{\centering natalia gajewska}
\date{}

% ------------- Creating a new block template ----------
\makeatletter
\def\th@newblock{%
  \normalfont 
  \def\inserttheoremblockenv{exampleblock}}
\theoremstyle{newblock}
\newtheorem{exercise}[theorem]{ćwiczenie}
\makeatother

\usepackage{environ}
\usepackage{verbatim}

\begin{document}
\begin{frame}
    \maketitle
\end{frame}
\begin{frame}[standout]
    zakładam że umiecie już tryb matematyczny bo jak nie to może być wam ciężej ale powinno być ok
\end{frame}
\begin{frame}[fragile]
    \frametitle{po co używać trybu chemicznego?}
    to jest bardzo uzasadnione pytanie! wszak wiele związków można napisać równie łatwo w trybie matematycznym i chemicznym, tu porównanie:
    \vspace{5mm}
        \begin{columns}
    \column{0.1\textwidth}
        \verb|\ce{H2O}|         
    \column{0.2\textwidth}
        \verb|$H_2O$|
    \end{columns}
    \vspace{5mm}
    \pause jednakże konwencją w chemii jest pisanie nazw związków bez kursywy, co różni się z konwencją matematyków, gdzie zmienne są pochylone. by to naprawić, tak naprawdę musielibyśmy pisać: \verb|$\mathrm{H_2O}$|.
\end{frame}

\begin{frame}[fragile]
    \frametitle{po co używać trybu chemicznego?}
    \ce{1/2 Ca(OH)2 + CO2 -> CaCO3 v + H2O} \\
    \vspace{5mm}
    tryb chemiczny: \\    
    \verb|\ce{1/2 Ca(OH)2 + CO2 -> CaCO3 v + H2O}| \\
    \vspace{5mm} tryb matematyczny: \\
    \verb|$\mathrm{\frac{1}{2} Ca(OH)_2 + CO_2| \\ \verb|\longrightarrow CaCO_3 \downarrow + H_2O}$|
\end{frame}

\begin{frame}
    pisanie jest po prostu sporo szybsze i prostsze \\
    tryb chemiczny daje nam bardzo miłe pisanie strzałek i innych symboli rzadkich w matematyce, a podstawowych w chemii
\end{frame}

\begin{frame}[fragile]
    warto jednak wspomnieć, że można i czasem* trzeba połączyć oba tryby by osiągnąć pewien efekt. tak jak w matematycznym, mamy możliwość pisania inline \verb|$ ... $| i blokowo \verb|\begin{equation} ... \end{equation}| \\
    niemniej, zazwyczaj wystarczy \verb|\ce{}| by osiągnąć swój cel. ta pisownia się "liczy" jako inline!
\end{frame}

\begin{frame}[standout]
    *o tym później ale proszę się nie bać
\end{frame}

\begin{frame}[fragile]
    \begin{exampleblock}{zaczynamy!}
        \verb|\usepackage[version=4]{mhchem}| \\
        \verb|\usepackage{amsmath}| \\
        \verb|\usepackage{chemfig}|
    \end{exampleblock}
\end{frame}

\begin{frame}[fragile]
    \frametitle{wzory sumaryczne i jony}
    \begin{exampleblock}{wzory sumaryczne}
        etanol, \ce{EtOH}, \verb|\ce{EtOH}| \\ \vspace{3mm}
        glukoza, \ce{C6H12O6}, \verb|\ce{C6H12O6}| \\ \vspace{3mm}
        wodorotlenek magnezu, \ce{Mg(OH)2}, \verb|\ce{Mg(OH)2}|
    \end{exampleblock}
\end{frame}

\begin{frame}[fragile]
    \frametitle{wzory sumaryczne i jony}
    \textbf{uważajcie na spacje!}
    \begin{exampleblock}{jony}
        anion sodu, \ce{Na+}, \verb|\ce{Na+}| \\ \vspace{1mm}
        sód plus (coś?), \ce{Na +}, \verb|\ce{Na +}| \\ \vspace{1mm}
        kation wapnia, \ce{Ca^{2+}}, \verb|\ce{Ca^{2+}}| \\ \vspace{1mm}
        \pause dwa wapnie bez elektronu? jednego?, \ce{Ca2+}, \verb|\ce{Ca2+}| \\ \vspace{3mm}
        \pause psst! tak samo piszemy stopnie utlenienia! \ce{Fe^{III}} to \verb|\ce{Fe^{III}}|
    \end{exampleblock}
\end{frame}

\newcounter{cwiczenia}
\newcommand\showmycounter{\addtocounter{cwiczenia}{1}\thecwiczenia}

\begin{frame}[fragile=singleslide]
    \frametitle{ćwiczenie \showmycounter \ - wzory sumaryczne}
    zapisz w trybie chemicznym 
    
    \begin{exampleblock}{związki}
        \begin{equation}
            \ce{CrO4^2-}
        \end{equation}
        \begin{equation}
            \ce{1/2Sb2O3}
        \end{equation}
        \begin{equation}
            \ce{[AgCl2]-}
        \end{equation}
    \end{exampleblock}
\end{frame}

\begin{frame}[fragile=singleslide]{rozwiązania!}
\begin{exampleblock}{związki}
    \begin{verbatim}
            \ce{CrO4^2-}

            \ce{1/2Sb2O3}

            \ce{[AgCl2]-}
    \end{verbatim}    
\end{exampleblock}
\end{frame}

\begin{frame}[fragile]
    aby napisać równania reakcji, przydadzą nam się strzałki! \vspace*{2mm}
    
    \ce{A -> B} to \verb|\ce{A -> B}|\\\vspace{1mm}
    \ce{A <- B} to \verb|\ce{A <- B}|\\\vspace{1mm}
    \ce{A <-> B} to \verb|\ce{A <-> B}|\\\vspace{1mm}
    \ce{A <--> B} to \verb|\ce{A <--> B}|\\\vspace{1mm}
    \ce{A <=> B} to \verb|\ce{A <=> B}|\\\vspace{1mm}
    \ce{A <=>> B} to \verb|\ce{A <=>> B}|\\\vspace{1mm}
    \ce{A <<=> B} to \verb|\ce{A <<=> B}|\\
    \vspace*{3mm}
    \pause przypominam że iupac nie lubi \ce{A <-> B} ale ja was nie uczę chemii tylko \LaTeX a
    \vspace{5mm}
    \pause    \begin{columns}
    \column{0.25\textwidth}
        \ce{ ^ } to \verb|\ce{ ^ }|         
    \column{0.25\textwidth}
        \ce{ v } to \verb|\ce{ v }|
    \end{columns}
\end{frame}
\begin{frame}[fragile]
    \ce{A ->[{góra}][{dól}] B} \vspace{5mm}
    \verb|\ce{A ->[{góra}][{dół}] B}|
\end{frame}
\begin{frame}[fragile=singleslide]
    \frametitle{ćwiczenie \showmycounter \ - równania reakcji} 
    \begin{exampleblock}{reakcje}
        \begin{equation}
            \ce{NaCl ->[{H_2O}] Na+ + Cl-}
        \end{equation}
        \begin{equation}
            \ce{^234_90Th -> ^0_-1\beta{} + ^234_91Pa}
        \end{equation}
        \begin{equation}
            \ce{2Al^0 + 3Cu^{II}(NO3)2 -> 2Al^{III}(NO3)3 + 3Cu^{0}}
        \end{equation}
    \end{exampleblock}
\end{frame}
\begin{frame}[fragile=singleslide]{rozwiązania!}
\begin{exampleblock}{reakcje}
    \begin{verbatim}
        \ce{NaCl ->[{H_2O}] Na+ + Cl-}

        \ce{^234_90Th -> ^0_-1\beta{} + ^234_91Pa}

        \ce{2Al^0 + 3Cu^{II}(NO3)2 ->
        2Al^{III}(NO3)3 + 3Cu^{0}}
    \end{verbatim} 
\end{exampleblock}
\end{frame}
\begin{frame}[fragile]
    \frametitle{rodniki i hydraty}
    \begin{verbatim}
        \ce{KCr(SO4)2*12H2O}
        \ce{KCr(SO4)2.12H2O}
        \ce{KCr(SO4)2 * 12 H2O} \end{verbatim}
     to wszystko daje \ce{KCr(SO4)2*12H2O}\vspace{2mm}
\end{frame}
\begin{frame}[fragile]
    \frametitle{hydraty i rodniki}
    wolne rodniki mają niesparowane elektrony, które oznaczamy kropką. dosłowną!\\ \vspace{2mm}
     \ce{OH^.} to \verb|\ce{OH^.}| \\
     \ce{CH3^.} to \verb|\ce{CH3^.}|
\end{frame}
\begin{frame}[fragile]
    ponadto, jak łatwo zauważyć, wiele znaków itp jest takich samych jak w trybie matematycznym:\\
    \ce{\alpha } to \verb|\ce{\alpha}| lub \verb|$\alpha$|\\\vspace{3mm}
    \pause i czasami konieczne jest ich łączenie!\\
    \ce{1/2}, czyli \verb|\ce{1/2}|, działa wyłącznie do prostych ułamków stechiometrycznych; do skomplikowanych użyjemy \verb|$\frac{licznik}{mianownik}$| jak zazwyczaj
\end{frame}
\begin{frame}[fragile=singleslide]
    \frametitle{ćwiczenie \showmycounter \ - misc} 
    \begin{exampleblock}{misc}
        \begin{equation}
            \ce{RNO2 <=>[+e] RNO2^{-.}}
        \end{equation}
        \begin{equation}
            \ce{CaCl2.2H2O}
        \end{equation}
        \begin{equation}
            K = \frac{[\ce{Hg^2+}][\ce{Hg}]}{[\ce{Hg2^2+}]}
        \end{equation}
    \end{exampleblock}
\end{frame}
\begin{frame}[fragile=singleslide]{rozwiązania!}
\begin{exampleblock}{misc}
    \begin{verbatim}
        \ce{RNO2 <=>[+e] RNO2^{-.}}

        \ce{CaCl2.2H2O}

        K = \frac{[\ce{Hg^2+}][\ce{Hg}]}
        {[\ce{Hg2^2+}]}
    \end{verbatim} 
\end{exampleblock}
\end{frame}
\begin{frame}
    \frametitle{myślnikowy problem}
    w trybie chemicznym możemy wykorzystać myślnik na kilka sposobów, jako:
    \begin{itemize}
        \item ładunek w jonie, choćby \ce{Cl-}
        \item minus w równaniu, jak \ce{3 - 2 = 1}
        \item pauza w nazwie, np. $\ce{$cis${-}[PtCl2(NH3)2]}$
        \item wiązanie pojedyncze w związku, np. \ce{C6H5-CHO}
    \end{itemize}
\end{frame}
\begin{frame}[fragile]
    \frametitle{myślnikowy problem}
    jako ładunek w jonie: \verb|\ce{Cl-}| \\ \vspace{1mm}
    minus w równaniu: \verb|\ce{3 - 2 = 1}| \\\vspace{1mm}
    pauza w nazwie: \verb|$\ce{$cis${-}[PtCl2(NH3)2]}$| \\\vspace{1mm}
    wiązanie pojedyncze: \verb|\ce{C6H5-CHO}|
    \\ \vspace{5mm}
\end{frame}
\begin{frame}[fragile]
    \frametitle{kursywa problem}
 tryb matematyczny, \verb|$...$|, wsadzony w \verb|\ce{...}|, podlega składni matematycznej, nie chemicznej - występuje kursywa etc. \\ chemiczny w matematycznym, \verb|$\ce{...}$|, ma składnię "chemiczną" \\ \vspace{3mm}
    \verb|\ce{$H2O$}| to będzie \ce{$H2O$}, a nie   $\ce{H2O}$.\\ \vspace{1mm}
    \verb|$\ce{$cis${-}buten}$| da nam $\ce{$cis${-}buten}$
\end{frame}
\begin{frame}[fragile]
    \frametitle{proste wzory strukturalne}
\verb|\ce{A-B=C#D}| daje nam 3 typy wiązań \ce{A-B=C#D} \\ \vspace{3mm}
\pause ale czasem hasztag nie działa jako potrójne wiązanie, i wtedy można napisać \verb|\bond{3}| między pierwiastkami. za pomocą \verb|\bond{coś}| można napisać też ciekawsze wiązania:\\\vspace{5mm}
\begin{columns}
    \column{0.25\textwidth}
    \ce{A\bond{->}B} \\\vspace{1mm}
    \ce{A\bond{...}B} \\\vspace{1mm}
    \ce{A\bond{~}B} \\\vspace{1mm}
    \ce{A\bond{~--}B}
    \column{0.35\textwidth}
    \verb|\ce{A\bond{->}B}| \\\vspace{1mm}
    \verb|\ce{A\bond{...}B}| \\\vspace{1mm}
    \verb|\ce{A\bond{~}B}| \\\vspace{1mm}
    \verb|\ce{A\bond{~--}B}|
\end{columns}
\end{frame}
\begin{frame}[standout]
    do bardziej zaawansowanych wzorów wykorzystamy chemfig. ale to za sekundkę
\end{frame}

\begin{frame}[fragile=singleslide]
    \frametitle{ćwiczenie \showmycounter \ - wiązania+} 
    \begin{exampleblock}{wiązania+}
        \begin{equation}
            \ce{O=C=O}
        \end{equation}
        \begin{equation}
            \ce{H-C#C-H}
        \end{equation}
        \begin{equation}
            \ce{(C5H5)2Fe2(\mu-CO)2(CO)2 <=> (C5H5)2Fe2(\mu-CO)2(CO)2}
        \end{equation}
        
    \end{exampleblock}
\end{frame}

\begin{frame}[fragile=singleslide]{rozwiązania!}
\begin{exampleblock}{wiązania+}
    \begin{verbatim}
        \ce{O=C=O}
        
        \ce{H-C#C-H}
        
        \ce{(C5H5)2Fe2(\mu-CO)2(CO)2 <=>
         (C5H5)2Fe2(\mu-CO)2(CO)2}
    \end{verbatim} 
\end{exampleblock}
\end{frame}
\begin{frame}
    \frametitle{brawo!}
    i tym samym skończyliśmy sekcję o mhchemie, która wystarczy wam do ok. 90\% waszych potrzeb chemicznych na sprawozdaniach. \\\vspace{5mm}
    \pause na reszcie tej sekcji porozmawiamy o pakiecie chemfig    
\end{frame}
\begin{frame}[standout]
    chemfig jest bazowany na TikZ, i zaczniemy od tego, co można nim zrobić, a nie jak działa
\end{frame}
\begin{frame}
    \centering nie chcę żebyście mi uciekli.
\end{frame}

\begin{frame}
    \frametitle{przykład}
    \begin{figure}
        \scalebox{1}{
        \chemfig{
            *6((=O)-N(-CH_3)-*5(-N=-N(-CH_3)-=)--(=O)-N(-H_3C)-)
            }
        }
  \caption{kofeina, 100\% \LaTeX, nwm czy etyczna}
  \label{kofeina}
    \end{figure}
i teraz mam ładną figurę, do której mogę się odwołać w tekście [\ref{kofeina}], elegancko podpisać, i ładnie skalować, bo to wektor ;>
\end{frame}
\begin{frame}[fragile]
    \frametitle{reveal kodu}
    \begin{verbatim}
    \begin{figure}
        \scalebox{1}{
        \chemfig{
            *6((=O)-N(-CH_3)-*5(-N=-N(-CH_3)-=)
            --(=O)-N(-H_3C)-)
            }
        }
  \caption{kofeina, 100\% \LaTeX, nwm czy etyczna}
  \label{kofeina}
    \end{figure}
    \end{verbatim}
\end{frame}
\begin{frame}[standout]
    jak widać, nie jest to takie straszne!
\end{frame}
\begin{frame}
    \begin{figure}
\scalebox{.65}{
\schemestart
\ce{O2^.-}\arrow (c1--c2){0}\ce{NO^.}\merge{v}(c1)(c2)--(c3)\ce{ONOO-}
\arrow (--c4){<=>[*{0}\parbox{1.2cm}{\ce{+ H+}\\ pK 6.6}][*{0}\ce{- H+}]}[-90,1.5]\ce{ONOOH}
\arrow(--c5){->[\ce{+ NO^.}][\ce{- HNO2}]}[0,1.5]\ce{NO2^.}
\arrow(--c6){->[\ce{+ NO^.}]}\ce{N2O3}
\arrow(@c4--c7){->[\ce{+ O2^.-}][\ce{- OH-},\ce{- O2}]}[180,1.5]\ce{NO2^.}
\arrow(@c7--c8){->[\ce{+ O2^.-}]}[180]\parbox{1.5cm}{\centering \ce{O2NOO-}\\or \\\ce{NO2- + O2}}
\arrow(@c3--@c7){->[\ce{+ CO2}][\ce{- CO3^.}]}[2.12]
\arrow(@c3--@c5){->[\ce{+ CO2}][\ce{- CO3^.-}]}[2.12]
\arrow(@c1--c9){->[\ce{+ O2^.-}, \ce{+ 2H+}][\ce{- O2}]}[180,1.5]\ce{H2O2}
\arrow(@c9--c10){->[\ce{Fe2+}][\ce{- OH-}]}[180]\ce{^.OH}
\arrow(@c2--c11){->[\ce{+ O2}]}$[$\ce{^.ONOO}$]$
\arrow(@c11--c12){->[\ce{+ ^.NO}]}\ce{2 NO2^.}
\schemestop
}
\end{figure}
\pause naturalnie nie wyłożę tego na tej sekcji bo byśmy stąd w życiu nie wyszli \\
\pause kodu też nie pokażę bo akurat ten jest zjebany
\end{frame}

\begin{frame}[fragile]
    chemfig służy do wzorów strukturalnych, i jak mogliście zauważyć, dla samych wzorów jest dosyć prosty i zwięzły (co oznacza że się go łatwo pisze a trudno czyta). 
\end{frame}

\begin{frame}[fragile]
    \frametitle{wiązania w chemfig}
    wiązania, tak jak w mchem, piszemy za pomocą - i =, ale mamy więcej ekscytujących opcji! \vspace{5mm}
\begin{columns}
    \column{0.33\textwidth}
        \centering
        \verb|\chemfig{A-B}| \\
        \chemfig{A-B}         
    \column{0.33\textwidth}
        \centering
        \verb|\chemfig{A=B}| \\
        \chemfig{A=B}
    \column{0.33\textwidth}
        \centering
        \verb|\chemfig{A~B}| \\
        \chemfig{A~B}
\end{columns}
\vspace{9mm}
\begin{columns}
    \column{0.33\textwidth}
        \centering
        \verb|\chemfig{A<B}| \\
        \chemfig{A<B}        
    \column{0.33\textwidth}
        \centering
        \verb|\chemfig{A<:B}| \\
        \chemfig{A<:B}
    \column{0.34\textwidth}
        \centering
        \verb|\chemfig{A<||\verb|B}|
        \chemfig{A<|B}
\end{columns}
\end{frame}

\begin{frame}[fragile]
    ale naturalnie, piszemy w chemfig bo chcemy wiele wiązań! \\ \vspace{3mm}
    \chemfig{A-B<C>:D~E}
    \vspace{3mm}
    \\ możemy ustalać ich kąty na kilka sposobów, z ogólną formą \verb|\chemfig{A-[kąt]B}|:
    \begin{itemize}
        \item z góry ustalone w pakiecie (czyli wielokrotności 45 stopni)
        \item bezwzględnie względem "kartki", od 0 do 360 stopni
        \item względnie do poprzedniego "atomu" - (właśnie, chemfig ma ich dziwną definicję) wtedy możemy wartości ujemne w prawo i dodatnie w lewo
        \item mamy też cool narzędzia do konkretnie związków o strukturze pierścieniowej
    \end{itemize}
\end{frame}

\begin{frame}[fragile]
\vspace{-2mm}
\begin{columns}
    \column{0.45\textwidth}
        \centering
        każde 1 to 45 stopni!
        \verb|{A-[1]B-[-1]C-[-2]D}| \\ \vspace{3mm}
        \chemfig{A-[1]B-[-1]C-[-2]D}       
    \column{0.45\textwidth}
        \centering
        po ":" w stopniach
        \verb|{A-[:45]B-[:-45]C-[:-90]D}| \\ \vspace{3mm}
        \chemfig{A-[:45]B-[:-45]C-[:-90]D} 
\end{columns}
\vspace{6mm}
\begin{columns}
    \column{0.01\textwidth}
    \column{0.98\textwidth}
        \centering
        po "::" mamy kąty względne do poprzedniego wiązania \\
        \verb|{A-[::45]B-[::-90]C-[::-45]D}| \\ \vspace{3mm}      
    \column{0.01\textwidth}
\end{columns}
\begin{columns}
    \column{0.45\textwidth}
        \centering
        \chemfig{A-[::45]B-[::-90]C-[::-45]D}       
    \column{0.45\textwidth}
        \centering
        \pause \chemfig{-[:45]-[:-45]-[:-90]} 
\end{columns}
\end{frame}
\begin{frame}[fragile=singleslide]
    \frametitle{ćwiczenie \showmycounter \ - proste wzory} 
    \begin{exampleblock}{kąty}
        \vspace{3mm}
        \begin{columns}
            \column{0.45\textwidth}
            \centering
            \chemfig{H-[:30]O-[:-30]H}
            \column{0.45\textwidth}
            \centering
            \chemfig{R-C-[::-60]O-[::-60]C-[::-60]R}
        \end{columns}
    \end{exampleblock}
\end{frame}

\begin{frame}[fragile=singleslide]{rozwiązania!}
     \begin{exampleblock}{kąty}
        \begin{columns}
    \column{0.01\textwidth}
    \column{0.98\textwidth}
        \centering
        \begin{verbatim}
            \chemfig{H-[:30]O-[:-30]H}

            \chemfig{R-C-[::-60]O-[::-60]C-[::-60]R}
        \end{verbatim}       
    \column{0.01\textwidth}
\end{columns}
\end{exampleblock}
\end{frame}

\begin{frame}[fragile]
    \frametitle{rozgałęzienia wzorów}
\vspace{10mm}
    \begin{columns}
    \column{0.45\textwidth}
        \centering
        \chemfig{A-B-C(-[:45]D)(-[:-45]E)} \\
        \vspace{3mm}
        \verb|A-B-C(-[:45]D)(-[:-45]E)}| \\    
    \column{0.45\textwidth}
        \centering
        \verb|{(-A)(-[2]B)(-[-2]D)(-[4]C)}| \\ \vspace{3mm}
        \chemfig{(-A)(-[2]B)(-[-2]D)(-[4]C)}       
\end{columns}    
\end{frame}

\begin{frame}[fragile=singleslide]
    \frametitle{ćwiczenie \showmycounter \ - rozgałęzienie} 
    \begin{exampleblock}{rozgałęzienia}
        \vspace{3mm}
        \begin{columns}
            \column{0.01\textwidth}
            \column{0.98\textwidth}
            \centering
            \chemfig{R-C(=[::+60]O)-[::-60]O-[::-60]C(=[::+60]O)-[::-60]R}
            \column{0.01\textwidth}
        \end{columns}
    \end{exampleblock}
\end{frame}

\begin{frame}[fragile=singleslide]{rozwiązania!}
     \begin{exampleblock}{rozgałęzienia}
        \begin{columns}
    \column{0.01\textwidth}
    \column{0.98\textwidth}
        \centering
        \begin{verbatim}
            \chemfig
            {R-C
            -[::-60]O
            -[::-60]C
            -[::-60]R}
        \end{verbatim}
        zamieniamy na: \\
        \begin{verbatim}
            \chemfig
            {R-C
            (=[::+60]O)
            -[::-60]O-[::-60]C
            (=[::+60]O)
            -[::-60]R}
        \end{verbatim}       
    \column{0.01\textwidth}
\end{columns}
\end{exampleblock}
\end{frame}

\begin{frame}[fragile]
    jeżeli mamy figurę opisaną kątami relatywnymi, czyli [::kąt], to całość możemy łatwo obrócić:\\ \vspace{3mm}
    \begin{verbatim}
    \chemfig{[:75]
    R-C(=[::+60]O)-[::-60]O-[::-60]C(=[::+60]O)-[::-60]R}
    \end{verbatim} \\ \vspace{7mm}
    \begin{columns}
    \column{0.01\textwidth}
    \column{0.98\textwidth}
    \centering
    \chemfig{[:75]R-C(=[::+60]O)-[::-60]O-[::-60]C(=[::+60]O)-[::-60]R}
    \column{0.01\textwidth}
    \end{columns}
\end{frame}

\begin{frame}[fragile]
    \frametitle{pierścienie}
    stworzenie jednego jest bardzo proste! \\ \vspace{6mm}
    \begin{columns}
    \column{0.45\textwidth}
        \centering
        \verb|\chemfig{*5(-----)}| \\ \vspace{2mm}
        \chemfig{*5(-----)}      
    \column{0.45\textwidth}
        \centering
        \verb|\chemfig{*6(-----)}| \\ \vspace{2mm}
        \chemfig{*6(------)}
    \end{columns}
\end{frame}

\begin{frame}[fragile]
    \frametitle{pierścienie}
    stworzenie jednego jest bardzo proste! \\ \vspace{6mm}
    \begin{columns}
    \column{0.45\textwidth}
        \centering
        \verb|\chemfig{*5(1-2-3-4-5-)}| \\ \vspace{2mm}
        \chemfig{*5(1-2-3-4-5-)}      
    \column{0.45\textwidth}
        \centering
        \verb|\chemfig{*6(1-2-3-4-5-6-)}| \\ \vspace{2mm}
        \chemfig{*6(1-2-3-4-5-6-)}
    \end{columns}
\end{frame}

\begin{frame}[fragile=singleslide]
    \frametitle{ćwiczenie \showmycounter \ - fenol} 
    \begin{exampleblock}{fenol}
        \vspace{3mm}
        \begin{columns}
            \column{0.01\textwidth}
            \column{0.98\textwidth}
            \centering
            \chemfig{*6(-=-=(-OH)-=)}
            \column{0.01\textwidth}
        \end{columns}
    \end{exampleblock}
\end{frame}
\begin{frame}[fragile=singleslide]{rozwiązania!}
     \begin{exampleblock}{rozgałęzienia}
        \begin{columns}
    \column{0.01\textwidth}
    \column{0.98\textwidth}
        \centering
        \begin{verbatim}
            \chemfig{*6(-=-=(-OH)-=)}
        \end{verbatim}
\end{columns}
\end{exampleblock}
btw, spróbujcie też \\ \verb|\chemfig{**6(----(-OH)--)}|
\end{frame}
\begin{frame}{co jeszcze robi chemfig?}
    \begin{itemize}
        \item czyści pierze bańki stawia
        \item można kolorować konkretne atomy i wiązania
        \item rozpisywać całe duże reakcje
        \item kopiować wzory z dokumentacji (bo jest ich dużo i często to łatwiejsze od pisania ab initio)
        \item definiować podjednostki molekuł (np. fosforany)
    \end{itemize}
\end{frame}
\begin{frame}[standout]
    dziękuję za uwagę!!
\end{frame}
\end{document}